\section*{ID7632: Small-Note: Bessel Functions (In Case You Actually Meant Bessel): J\_n(·)}
\paragraph{Metadata.}
PiID: 0326591840443 | RTEID: RTE:P29636.7509:T4.8500 | UUID: 239a366d-ef3c-5ffd-baee-bae44cd3c0fe

\paragraph{Canonical Statement.}
- A sinusoidally-modulated phase process produces sidebands whose amplitudes are Bessel functions \textbackslash{}(J\_n(\textbackslash{}cdot)\textbackslash{}). - If creation depends on modulation depth (e.g., due to a periodic drive), the creation probability can be written with coefficients \textbackslash{}(J\_n\textasciicircum{}2(\textbackslash{}kappa)\textbackslash{}) for harmonic \textbackslash{}(n\textbackslash{}). - So if you meant *Bessel* rather than *beat*, the creation kernel would weight harmonics by Bessel amplitudes and could naturally explain higher-harmonic spin–orbit resonances (like Mercury 3:2).

\paragraph{Canonical Equation.}
\[
J_n(\cdot)
\]

\paragraph{Dictionary.}
Relation: J\_n(·)

\paragraph{Encyclopedia.}
Small-Note: Bessel Functions (In Case You Actually Meant Bessel): J\_n(·) is a symbolic expression, written J\_n(·). It introduces a symbol or relation used elsewhere in the framework. It is built from J\_n. Within the Trawin Zero Point registry it serves as one element of the framework's mathematical narrative.
